\chapter{Introduction}

\section{Contexte industriel}
La disponibilité des équipements est un facteur déterminant de la compétitivité
industrielle. Dans une usine, une heure d'arrêt d'une ligne de production peut coûter
plusieurs milliers d'euros, et les défaillances de machines tournantes (pompes,
ventilateurs, compresseurs, réducteurs, turbines) figurent parmi les causes les plus
fréquentes d'indisponibilité. Or, plus de la moitié de ces défaillances sont d'origine
mécanique et se manifestent, des semaines voire des mois à l'avance, par une évolution
caractéristique du comportement vibratoire de la machine.

On distingue traditionnellement trois stratégies de maintenance :
\begin{itemize}[leftmargin=1.5em]
  \item la \textbf{maintenance corrective}, qui intervient après la panne. Simple à
        organiser, elle est la plus coûteuse : arrêt subi, dégâts collatéraux, perte de
        production, risque sécurité.
  \item la \textbf{maintenance préventive systématique}, qui intervient à intervalles
        fixes (calendaires ou à compteur). Elle réduit les pannes mais engendre un
        gaspillage : on remplace des composants encore sains, et l'on n'élimine pas les
        défaillances précoces survenant entre deux interventions.
  \item la \textbf{maintenance prédictive} (ou conditionnelle, \emph{Condition-Based
        Maintenance}, CBM), qui intervient \emph{au bon moment}, en fonction de l'état
        réel de la machine mesuré en continu. C'est la stratégie la plus efficiente, mais
        aussi la plus exigeante techniquement.
\end{itemize}

La maintenance prédictive repose sur un principe simple : une machine qui se dégrade
\og parle \fg{} à travers ses signaux physiques (vibration, température, courant,
acoustique). L'enjeu est de \emph{savoir l'écouter} : capter ces signaux, en extraire des
indicateurs pertinents, reconnaître les signatures de défaut, et estimer le temps restant
avant la défaillance.

\section{Problématique}
La mise en œuvre d'une solution de maintenance prédictive crédible se heurte à plusieurs
difficultés concrètes :
\begin{enumerate}[leftmargin=1.6em]
  \item \textbf{Détecter tôt} une dégradation, alors que ses premières manifestations sont
        noyées dans le bruit et masquées par les conditions d'exploitation variables.
  \item \textbf{Diagnostiquer précisément} le type de défaut (balourd, désalignement,
        défaut de roulement, jeu mécanique, défaut d'engrenage) pour orienter
        l'intervention.
  \item \textbf{Quantifier} l'état de santé et la durée de vie restante de façon
        exploitable pour la planification.
  \item \textbf{Maîtriser le taux de fausses alarmes} : un système qui alerte à tort
        perd la confiance des équipes et finit désactivé.
  \item \textbf{Rendre la solution actionnable et traçable} : transformer une alerte en
        intervention planifiée, et pouvoir justifier chaque décision a posteriori.
\end{enumerate}

À ces difficultés techniques s'ajoute une contrainte de \emph{crédibilité industrielle} :
les indicateurs et diagnostics produits doivent parler le langage normalisé des ingénieurs
de fiabilité (zones de sévérité ISO, MTBF, MTTR, disponibilité) plutôt que des métriques
propriétaires non vérifiables.

\section{Objectifs du projet}
L'objectif est de concevoir et de réaliser une plateforme logicielle complète, nommée
\textbf{PrognoSense}, qui couvre l'ensemble de la chaîne de maintenance prédictive :
\begin{itemize}[leftmargin=1.5em]
  \item acquérir et traiter des signaux vibratoires issus de plusieurs natures de machines ;
  \item extraire des indicateurs physiquement interprétables et diagnostiquer les défauts ;
  \item entraîner, comparer et sélectionner des modèles d'apprentissage automatique de
        classification de défauts et de pronostic de durée de vie ;
  \item détecter des anomalies inconnues par apprentissage non supervisé ;
  \item évaluer la sévérité selon la norme ISO~10816/20816 ;
  \item synthétiser un \emph{indice de santé} et suivre la flotte en temps réel ;
  \item transformer les alertes en ordres de travail transmissibles à une GMAO ;
  \item mesurer la fiabilité réelle (MTBF, MTTR, disponibilité, ROI, fausses alarmes) ;
  \item garantir la confiance : explicabilité, calibration, audit, versioning des modèles ;
  \item exposer le tout via une interface web ergonomique et un copilot conversationnel.
\end{itemize}

\section{Démarche et organisation du rapport}
La démarche a suivi une progression incrémentale : maîtrise du pipeline de traitement du
signal et des données, entraînement et benchmark des modèles, construction de la plateforme
applicative (backend/frontend), puis industrialisation (normes, connectivité, boucle
fermée, confiance).

Le rapport est organisé comme suit. Le chapitre~\ref{chap:etatart} présente l'état de l'art
de la maintenance prédictive et le positionnement vis-à-vis des solutions du marché. Le
chapitre~\ref{chap:datasets} analyse les jeux de données utilisés. Le
chapitre~\ref{chap:archi} décrit l'architecture générale. Le chapitre~\ref{chap:pipeline}
détaille le pipeline de traitement du signal et l'extraction des indicateurs. Le
chapitre~\ref{chap:bench} présente le benchmark des cinq modèles. Le
chapitre~\ref{chap:anomalie} traite la détection d'anomalies. Les
chapitres~\ref{chap:backend} et~\ref{chap:frontend} décrivent respectivement le backend et
le frontend. Le chapitre~\ref{chap:indus} couvre l'industrialisation (ISO, connectivité,
GMAO, KPIs réels, copilot, confiance). Le chapitre~\ref{chap:defis} relate les défis
techniques rencontrés. Le chapitre~\ref{chap:resultats} synthétise et interprète les
résultats. Le chapitre~\ref{chap:limites} discute les limites et perspectives, et le
chapitre~\ref{chap:conclusion} conclut.
